\chapter{Related Work}

\section{Feature Engineering and Statistical Models}

The earliest wave of legal outcome prediction relied on feature engineering and classical statistical learning, representing cases using bag-of-words, TF--IDF vectors, or manually designed legal features, and training linear models, SVMs, or tree-based classifiers. Aletras et al.\ showed that n-gram and topic features extracted from ECHR case text could predict violation judgments with around 79\% accuracy, with case facts emerging as the strongest predictive signal~\cite{aletras2016echr}. In the US context, Katz et al.\ built a time-evolving random forest over institutional metadata and historical voting patterns, achieving over 70\% accuracy predicting Supreme Court outcomes across nearly two centuries~\cite{katz2017supreme}. Medvedeva et al.\ extended ECHR prediction work and demonstrated that predicting genuinely future cases from past ones substantially lowers performance, raising early concerns about temporal leakage~\cite{medvedeva2020echr}.

These approaches remain relevant as strong, truncation-free baselines: TF--IDF classifiers can be competitive on long documents precisely because they avoid the input-length constraints that trouble neural models~\cite{mamakas2022longlegal}. Their central limitation, however, is that flat lexical representations cannot capture discourse order, typed legal relations, or cross-case structure. A broader methodological critique by Medvedeva et al.\ found that the majority of LJP publications do not perform genuine prospective prediction, instead relying on unrealistic inputs or temporally inconsistent splits~\cite{medvedeva2023doitright}. This critique is directly relevant here: the present thesis treats leakage control as a core methodological contribution.

\section{Text-Based Deep Learning}

The second major wave replaced manual features with learned neural representations. Domain-adapted transformers quickly became dominant; LEGAL-BERT demonstrated that pre-training on legal corpora---legislation, court cases, and contracts---yields consistent gains over generic BERT on legal classification tasks~\cite{chalkidis2020legalbert}. The LexGLUE benchmark consolidated several legal NLU datasets and confirmed that legal-domain models systematically outperform generic ones across diverse tasks~\cite{chalkidis2022lexglue}.

The central weakness of standard transformers in law is document length. Most judgments far exceed the 512-token window of vanilla BERT, forcing researchers to truncate, use sliding windows, or adopt hierarchical architectures. Work on modified Longformer architectures shows that extending context to 8,192 sub-words improves performance on long legal classification tasks, but at substantial computational cost, and even then hierarchical TF--IDF variants remain competitive in some settings~\cite{mamakas2022longlegal}. Hierarchical encoders are attractive because they mirror legal intuition---facts, arguments, and operative text play distinct roles---but they remain document-centric: they model internal case structure while ignoring cross-case links through shared statutes or precedents. The present thesis departs from purely text-based work by treating text embeddings as node features within a larger relational representation, rather than as the complete model.

\section{Graph-Based Approaches}

Graph-based methods were developed precisely to model the relational structure that flat text pipelines miss. TopJudge formalised dependencies among legal subtasks---law articles, charges, and penalties---as a directed acyclic graph, showing that joint prediction over structured outputs improves performance compared to treating each subtask independently~\cite{zhong2018topjudge}. Subsequent work expanded this relational view by incorporating retrieved precedents and law-article graphs, with more recent systems using LLM-and-retrieval pipelines to bring external legal scaffolding into prediction~\cite{wu2023pljp, santosh2024precedents}.

Within Indian LJP, graph-based work is growing. Khatri et al.\ explored GNNs for Indian legal judgment prediction and raised important questions about the role of case-graph design and demographic shortcuts~\cite{khatri2023indianGNN}. More recent large-scale Indian work has extended coverage and explored RAG-style pipelines~\cite{nigam2025nyayaanumana}. Existing graph approaches nonetheless leave gaps that this thesis addresses: many are built for civil-law charge-prediction datasets, rely on a single court, or use graph structure to improve prediction without comparable attention to decision-text or identity leakage. The present work is closest in spirit to this line but narrows its claim: the graph should improve prediction by exposing legal structure across cases while explicitly blocking shortcuts that rely on post-decision information.

\section{Heterogeneous GNNs and the Motivation for HGT}

When a legal problem is represented as a graph, the choice of GNN architecture matters because legal graphs contain multiple semantically distinct node and edge types. A homogeneous GNN is too blunt: a case node, a statute node, and a precedent node play different epistemic roles, and edges such as \emph{cites\_precedent} and \emph{heard\_in} should not be collapsed into a single message-passing rule.

R-GCN introduced relation-specific weight matrices for multi-relational graphs and is well-suited to knowledge graph completion tasks~\cite{schlichtkrull2018rgcn}. Its limitation for the present thesis is that relation-specific parameterisation can become coarse as the number of types grows, and it does not model attention over heterogeneous neighbourhoods. HAN addressed heterogeneity through hierarchical attention over meta-paths, learning node-level and semantic-level importance separately~\cite{wang2019han}. However, it requires manually specified meta-paths, which hard-codes assumptions that may not transfer across legal domains. HGT provides a more flexible alternative: it uses node-type-dependent projections and edge-type-dependent attention, allowing the model to learn how much weight to assign to different heterogeneous neighbours without requiring handcrafted meta-path templates~\cite{hu2020hgt}. For the legal graph constructed in this thesis---which includes many node types, many edge types, and a deliberately cross-domain structure---HGT fits the design requirements better than R-GCN or HAN.

\section{Indian Legal NLP and OpenNyAI}

Indian legal NLP developed later than English-language benchmarks for US or European law, in part due to data and tooling scarcity. Early work on rhetorical role identification showed that Indian judgments have a recoverable discourse structure---preamble, facts, arguments, ratio, and ruling---that is useful for downstream tasks but not explicitly marked in raw text~\cite{bhattacharya2019rhetoricalroles}. This made legal document structuring a first-class problem rather than a preprocessing convenience.

OpenNyAI is the most practically important development for this thesis, bringing together legal named entity recognition, rhetorical-role prediction, and summarisation for Indian court judgments into an open, modular pipeline~\cite{kalamkar2022ner, opennyai2024docs, opennyai2024about}. On the benchmarking side, ILDC introduced a large corpus of Indian Supreme Court cases for judgment prediction and explanation~\cite{malik2021ildc}, while IL-TUR broadened the scope to a multi-task legal understanding and reasoning benchmark~\cite{joshi2024iltur}. Newer rhetorical-role resources such as LegalSeg confirm that structural modelling continues to matter for Indian judgments~\cite{nigam2025legalseg}.

These contributions are substantial, but they do not resolve the core challenge addressed here: how to represent Indian cases so that a model can exploit legal structure across cases without quietly learning identity or post-decision shortcuts. This thesis builds on the Indian legal NLP ecosystem rather than apart from it, combining OpenNyAI's pipeline components with a leakage-audited heterogeneous graph.

\section{Gaps This Work Addresses}

The literature reviewed above shows that legal judgment prediction has advanced substantially, but along partly disconnected lines. Classical models demonstrated that outcomes are statistically predictable; transformer models improved semantic representation; graph methods showed the value of relational structure; Indian NLP resources made domain-specific preprocessing feasible. Yet most prior work sacrifices at least one of the following: \emph{scale}, \emph{cross-domain coverage}, \emph{cross-case legal structure}, or \emph{leakage-aware experimental design}.

What remains comparatively rare is an Indian LJP system that operates over a large multi-domain corpus, consolidates multiple hearings per case, extracts legal entities and rhetorical structure, builds an explicitly heterogeneous cross-case graph, and evaluates under a design intended to resist both decision-text leakage and identity-based shortcuts. That is the niche this thesis occupies. Its claim is narrow and, for that reason, defensible: when Indian court judgments are converted into a leakage-audited heterogeneous graph, cross-case legal structure becomes usable in ways that flat text pipelines cannot easily match.